\chapter{Lecture, débogage et conception logicielle}\label{ch:iv3:reading-debugging-and-software-design}

\begin{ivtrains}
Lire un court listing et dire ce qu'il fait réellement, puis nommer le bogue, l'outil qui permettrait de le prouver et le plus petit correctif. Les premières questions sont celles que l'on attend de repérer en quelques secondes. Les suivantes demandent comment chasser une faute qu'on ne voit pas, comment tester un morceau de code de trading et comment le structurer pour que la même logique tourne en direct et en backtest. Quelques énoncés sont en Java ou en C\#, car les banques les posent dans ces langages.
\end{ivtrains}

\begin{question}[$\star$ \iqroles{QD, Dev, LLE} \iqfirm{toutes} \iqlevel{junior}]\label{iq:iv3:reading-debugging-and-software-design:1}
Compilé avec g++ 11.4 sur x86-64 Linux, qu'affiche le \cref{lst:iv3:rd:unsigned}~? Y a-t-il un comportement indéfini ici~? Donnez le plus petit correctif pour la comparaison.
\omcode{code/interviews-3/09-reading-debugging-and-software-design/cpp/unsigned_diff.cpp}{4}{13}{Comparaison de deux tailles.\protect\hfil\protect\break}\label{lst:iv3:rd:unsigned}
\end{question}

\begin{question}[$\star$ \iqroles{QR, QT, QD, Dev, DS} \iqfirm{toutes} \iqlevel{junior}]\label{iq:iv3:reading-debugging-and-software-design:2}
Qu'affiche le \cref{lst:iv3:rd:late}, et pourquoi~? Faites en sorte qu'il affiche \texttt{[10, 20, 30]}. Une réécriture de la boucle en compréhension de liste aiderait-elle~?
\omcode{code/interviews-3/09-reading-debugging-and-software-design/python/rd_late_binding.py}{2}{9}{Trois fonctions de mise à l'échelle.\protect\hfil\protect\break}\label{lst:iv3:rd:late}
\end{question}

\begin{question}[$\star$ \iqroles{QD, Dev, LLE} \iqfirm{toutes} \iqlevel{junior}]\label{iq:iv3:reading-debugging-and-software-design:3}
Le \Cref{lst:iv3:rd:erase} est censé supprimer chaque quantité nulle. Compilé avec g++ 11.4, \texttt{-std=c++20 -O2}, il affiche quelque chose. Quoi, et pourquoi~? Que se passe-t-il si le dernier élément est le zéro~? Quel outil met la faute en évidence, et quel est le correctif~?
\omcode{code/interviews-3/09-reading-debugging-and-software-design/cpp/erase_loop.cpp}{4}{13}{Suppression des zéros dans une boucle.\protect\hfil\protect\break}\label{lst:iv3:rd:erase}
\end{question}

\begin{question}[$\star$ \iqroles{QD, Dev, LLE} \iqfirm{toutes} \iqlevel{junior}]\label{iq:iv3:reading-debugging-and-software-design:4}
Qu'affiche le \cref{lst:iv3:rd:mapindex}~? Il y a deux fautes dans \texttt{allowed}. L'une est un mauvais résultat et l'autre un effet de bord. Quelle est la nature de chacune, et comment feriez-vous pour que le compilateur détecte la seconde~?
\omcode{code/interviews-3/09-reading-debugging-and-software-design/cpp/map_index.cpp}{5}{15}{Une vérification de limite.\protect\hfil\protect\break}\label{lst:iv3:rd:mapindex}
\end{question}

\begin{question}[$\star$ \iqroles{QD, Dev, LLE} \iqfirm{toutes} \iqlevel{junior}]\label{iq:iv3:reading-debugging-and-software-design:5}
Qu'affiche le \cref{lst:iv3:rd:ctor}, et pourquoi~? Supposez que \texttt{name} soit pure virtuelle dans \texttt{Strategy}, et que le constructeur l'appelle par l'intermédiaire d'une autre fonction membre. Que se passe-t-il alors~? Comment obtiendriez-vous le comportement voulu par son auteur~?
\omcode{code/interviews-3/09-reading-debugging-and-software-design/cpp/ctor_virtual.cpp}{3}{16}{Un appel virtuel pendant la construction.\protect\hfil\protect\break}\label{lst:iv3:rd:ctor}
\end{question}

\begin{question}[$\star$ \iqroles{QD, Dev, LLE} \iqfirm{banque} \iqlevel{junior}]\label{iq:iv3:reading-debugging-and-software-design:6}
Quels sont les trois nombres qu'affiche le \cref{lst:iv3:rd:slicing}~? Nommez l'effet, donnez le correctif usuel pour ce livre, et dites comment faire en sorte que l'erreur fasse échouer la compilation.
\omcode{code/interviews-3/09-reading-debugging-and-software-design/cpp/slicing.cpp}{4}{22}{Un livre d'instruments.\protect\hfil\protect\break}\label{lst:iv3:rd:slicing}
\end{question}

\begin{question}[$\star$ \iqroles{QD, Dev, LLE} \iqfirm{crypto, teneur de marché} \iqlevel{junior}]\label{iq:iv3:reading-debugging-and-software-design:7}
Avec rustc 1.97 sur x86-64, que fait le \cref{lst:iv3:rd:usize} sous \texttt{cargo run}~? Et sous \texttt{cargo run -{}-release}~? L'un de ces deux cas est-il unsafe au sens de Rust~? Corrigez \texttt{imbalance} de deux façons, selon ce que veut l'appelant.
\omcode{code/interviews-3/09-reading-debugging-and-software-design/rust/src/bin/imbalance.rs}{1}{9}{Un déséquilibre.\protect\hfil\protect\break}\label{lst:iv3:rd:usize}
\end{question}

\begin{question}[$\star$ \iqroles{QR, QT, QD, Dev, DS} \iqfirm{toutes} \iqlevel{junior}]\label{iq:iv3:reading-debugging-and-software-design:8}
Quelles sont les deux lignes qu'affiche le \cref{lst:iv3:rd:floatkeys} sous CPython~? Un collègue stocke des niveaux de prix dans un dictionnaire indexé par des flottants, et des exécutions disparaissent de temps en temps. Que lui répondez-vous~?
\omcode{code/interviews-3/09-reading-debugging-and-software-design/python/rd_float_keys.py}{1}{4}{Flottants utilisés comme clés, et arrondis.\protect\hfil\protect\break}\label{lst:iv3:rd:floatkeys}
\end{question}

\begin{question}[$\star$ \iqroles{QD, Dev} \iqfirm{banque} \iqlevel{junior}]\label{iq:iv3:reading-debugging-and-software-design:9}
Sous OpenJDK 11 avec les réglages par défaut, qu'affiche le \cref{lst:iv3:rd:boxes}~? Lesquels de ces résultats le langage Java garantit-il, et lesquels relèvent de cette JVM~?
\omcode{code/interviews-3/09-reading-debugging-and-software-design/java/Boxes.java}{1}{9}{Comparaison d'objets enveloppes et de chaînes.\protect\hfil\protect\break}\label{lst:iv3:rd:boxes}
\end{question}

\begin{question}[$\star\star$ \iqroles{QD, Dev, LLE} \iqfirm{toutes} \iqlevel{junior, confirmé}]\label{iq:iv3:reading-debugging-and-software-design:10}
Le \Cref{lst:iv3:rd:dtor} se compile sans avertissement avec g++ 11.4, \texttt{-std=c++20 -O2 -Wall -Wextra -Werror}. Qu'affiche-t-il~? Le programme est-il correct~? Que dirait AddressSanitizer~? La réponse changerait-elle avec \texttt{std::make\_shared}~?
\omcode{code/interviews-3/09-reading-debugging-and-software-design/cpp/no_virtual_dtor.cpp}{5}{19}{Un ordre derrière un pointeur de base.\protect\hfil\protect\break}\label{lst:iv3:rd:dtor}
\end{question}

\begin{question}[$\star\star$ \iqroles{QD, Dev, LLE} \iqfirm{toutes} \iqlevel{junior, confirmé}]\label{iq:iv3:reading-debugging-and-software-design:11}
Le \Cref{lst:iv3:rd:exleak} affiche le total attendu alors qu'il tourne depuis des mois. Qu'est-ce qui cloche~? Combien perd-il, et sur quel chemin~? Quel outil le révèle, et quel est le correctif~? Votre collègue propose un \texttt{catch}, un \texttt{delete} et une nouvelle levée de l'exception dans \texttt{handle}. L'acceptez-vous~?
\omcode{code/interviews-3/09-reading-debugging-and-software-design/cpp/exception_leak.cpp}{5}{29}{Traitement d'un ordre.\protect\hfil\protect\break}\label{lst:iv3:rd:exleak}
\end{question}

\begin{question}[$\star\star$ \iqroles{QD, Dev, LLE} \iqfirm{toutes} \iqlevel{junior, confirmé}]\label{iq:iv3:reading-debugging-and-software-design:12}
Le \Cref{lst:iv3:rd:temprange} est compilé en C++20 avec g++ 11.4. Est-il défini~? Rendez maintenant \texttt{levels\_} public et parcourez \texttt{latest().levels\_}. Est-ce défini~? Qu'a changé le C++23~? Donnez un correctif pour le C++20 et un moyen de faire rejeter l'erreur au compilateur.
\omcode{code/interviews-3/09-reading-debugging-and-software-design/cpp/temp_range.cpp}{4}{19}{Parcours du dernier instantané.\protect\hfil\protect\break}\label{lst:iv3:rd:temprange}
\end{question}

\begin{question}[$\star\star$ \iqroles{QD, Dev, LLE} \iqfirm{toutes} \iqlevel{junior, confirmé}]\label{iq:iv3:reading-debugging-and-software-design:13}
Deux threads exécutent le \cref{lst:iv3:rd:guard}. Le code compile avec g++ 11.4 sous \texttt{-Wall -Wextra -Werror}, et l'auteur attend 200000. Qu'est-ce qui ne va pas~? Montrez le correctif. Que se passe-t-il si vous écrivez la même ligne avec des parenthèses au lieu des accolades~?
\omcode{code/interviews-3/09-reading-debugging-and-software-design/cpp/guard_temp.cpp}{5}{19}{Un compteur derrière un mutex.\protect\hfil\protect\break}\label{lst:iv3:rd:guard}
\end{question}

\begin{question}[$\star\star$ \iqroles{QD, Dev, LLE} \iqfirm{toutes} \iqlevel{junior, confirmé}]\label{iq:iv3:reading-debugging-and-software-design:14}
Le \Cref{lst:iv3:rd:lockorder} affiche \texttt{1000 1000} à chaque fois, car les threads ne se chevauchent jamais. Quel est le bogue dans \texttt{transfer}~? ThreadSanitizer le signalerait-il sur cette exécution~? Corrigez-le. Puis~: que fait votre correctif lors d'un transfert d'un compte vers lui-même~?
\omcode{code/interviews-3/09-reading-debugging-and-software-design/cpp/lock_order.cpp}{5}{21}{Un transfert entre comptes.\protect\hfil\protect\break}\label{lst:iv3:rd:lockorder}
\end{question}

\begin{question}[$\star\star$ \iqroles{QD, Dev, LLE} \iqfirm{toutes} \iqlevel{junior, confirmé}]\label{iq:iv3:reading-debugging-and-software-design:15}
Trouvez le bogue dans \texttt{pop} dans le \cref{lst:iv3:rd:cvif} et dites exactement comment il se produit avec ces deux consommateurs. Comment rendriez-vous la défaillance visible~? Corrigez-le. Le remplacement de \texttt{notify\_all} par \texttt{notify\_one} suffirait-il aussi à le corriger~?
\omcode{code/interviews-3/09-reading-debugging-and-software-design/cpp/cv_if.cpp}{8}{37}{Une file à deux consommateurs.\protect\hfil\protect\break}\label{lst:iv3:rd:cvif}
\end{question}

\begin{question}[$\star\star$ \iqroles{QD, Dev, LLE} \iqfirm{toutes} \iqlevel{junior, confirmé}]\label{iq:iv3:reading-debugging-and-software-design:16}
Qu'affiche le \cref{lst:iv3:rd:constmove}~? Son auteur était persuadé qu'il n'y aurait aucune copie. Pourquoi y en a-t-il une~? Expliquez ensuite pourquoi un constructeur de déplacement sans \texttt{noexcept} peut aussi réintroduire des copies, à l'intérieur de \texttt{std::vector}.
\omcode{code/interviews-3/09-reading-debugging-and-software-design/cpp/const_move.cpp}{5}{25}{Déplacement vers un membre.\protect\hfil\protect\break}\label{lst:iv3:rd:constmove}
\end{question}

\begin{question}[$\star\star$ \iqroles{QD, Dev, LLE} \iqfirm{teneur de marché, hedge fund} \iqlevel{junior, confirmé}]\label{iq:iv3:reading-debugging-and-software-design:17}
Le \Cref{lst:iv3:rd:callback} affiche \texttt{1} puis \texttt{done} lorsqu'il est compilé normalement. Est-il correct~? Que signalerait AddressSanitizer, et à quelle ligne~? Concevez le correctif. Qu'exige-t-il de la classe pour que personne ne puisse la casser par mégarde~?
\omcode{code/interviews-3/09-reading-debugging-and-software-design/cpp/callback_this.cpp}{7}{33}{Une stratégie s'abonne à un flux de données.\protect\hfil\protect\break}\label{lst:iv3:rd:callback}
\end{question}

\begin{question}[$\star\star$ \iqroles{QD, Dev, DS} \iqfirm{toutes} \iqlevel{junior, confirmé}]\label{iq:iv3:reading-debugging-and-software-design:18}
Qu'affiche la dernière ligne du \cref{lst:iv3:rd:hash} sous CPython, et pourquoi~? Sur quoi cette classe devrait-elle calculer son hachage~? Que se passe-t-il si vous supprimez simplement son \texttt{\_\_hash\_\_}~?
\omcode{code/interviews-3/09-reading-debugging-and-software-design/python/rd_mutable_hash.py}{1}{16}{Des ordres dans un ensemble.\protect\hfil\protect\break}\label{lst:iv3:rd:hash}
\end{question}

\begin{question}[$\star\star$ \iqroles{QR, QT, QD, Dev, DS} \iqfirm{toutes} \iqlevel{junior, confirmé}]\label{iq:iv3:reading-debugging-and-software-design:19}
Le \Cref{lst:iv3:rd:cache} est un pricer qui met ses résultats en cache. Qu'affichent les deux lignes~? Un processus de longue durée crée un \texttt{Pricer} par courbe, et sa mémoire ne cesse de croître. Pourquoi~? Donnez un correctif qui conserve le cache.
\omcode{code/interviews-3/09-reading-debugging-and-software-design/python/rd_method_cache.py}{2}{21}{Un pricer avec cache.\protect\hfil\protect\break}\label{lst:iv3:rd:cache}
\end{question}

\begin{question}[$\star\star$ \iqroles{QR, QT, QD, Dev, DS} \iqfirm{toutes} \iqlevel{junior, confirmé}]\label{iq:iv3:reading-debugging-and-software-design:20}
Le \Cref{lst:iv3:rd:dst} mesure une journée à New York sous Python 3.10. Qu'affichent les deux lignes~? Laquelle correspond à la durée de cette journée~? Citez un calcul, dans une société de trading, qui donnerait silencieusement un résultat faux avec la première.
\omcode{code/interviews-3/09-reading-debugging-and-software-design/python/rd_dst.py}{1}{8}{La durée d'une journée.\protect\hfil\protect\break}\label{lst:iv3:rd:dst}
\end{question}

\begin{question}[$\star\star$ \iqroles{QD, Dev, LLE} \iqfirm{crypto, teneur de marché} \iqlevel{junior, confirmé}]\label{iq:iv3:reading-debugging-and-software-design:21}
Qu'affiche le \cref{lst:iv3:rd:rccycle}, avec rustc 1.97~? Quels messages ne voit-on jamais, et pourquoi~? Est-ce un bogue selon les propres règles de Rust~? Corrigez-le en changeant un seul type, et dites ce qu'affiche la version corrigée.
\omcode{code/interviews-3/09-reading-debugging-and-software-design/rust/src/bin/rc_cycle.rs}{1}{20}{Deux nœuds.\protect\hfil\protect\break}\label{lst:iv3:rd:rccycle}
\end{question}

\begin{question}[$\star\star$ \iqroles{QD, Dev, LLE} \iqfirm{crypto, teneur de marché} \iqlevel{confirmé}]\label{iq:iv3:reading-debugging-and-software-design:22}
\Cref{lst:iv3:rd:guardmatch} se compile sans avertissement avec rustc 1.97. Que se passe-t-il à l'exécution, et pourquoi~? L'édition 2024 change-t-elle quelque chose~? Corrigez-le en ajoutant une seule ligne.
\omcode{code/interviews-3/09-reading-debugging-and-software-design/rust/src/bin/guard_match.rs}{1}{13}{Remettre un élément en file.\protect\hfil\protect\break}\label{lst:iv3:rd:guardmatch}
\end{question}

\begin{question}[$\star\star$ \iqroles{QD, Dev} \iqfirm{banque} \iqlevel{junior, confirmé}]\label{iq:iv3:reading-debugging-and-software-design:23}
Sous OpenJDK 11, qu'affiche le \cref{lst:iv3:rd:equalsonly}~? Pourquoi la deuxième ligne surprend-elle, étant donné la première~? Énoncez la règle qui a été enfreinte et la correction. Donnez ensuite une autre façon de casser une clé, même lorsque les deux méthodes sont écrites.
\omcode{code/interviews-3/09-reading-debugging-and-software-design/java/EqualsOnly.java}{1}{21}{Une classe de clé.\protect\hfil\protect\break}\label{lst:iv3:rd:equalsonly}
\end{question}

\begin{question}[$\star\star$ \iqroles{QD, Dev} \iqfirm{banque} \iqlevel{junior, confirmé}]\label{iq:iv3:reading-debugging-and-software-design:24}
Sous OpenJDK 11, \cref{lst:iv3:rd:lostmerge} devrait afficher 400000. Ce n'est pas le cas. La table est un \texttt{ConcurrentHashMap}~: qu'est-ce qui a mal tourné~? Faites tenir la correction en une seule ligne. Qu'utiliseriez-vous si le compteur était mis à jour depuis un très grand nombre de threads et lu rarement~?
\omcode{code/interviews-3/09-reading-debugging-and-software-design/java/LostMerge.java}{1}{22}{Ajouter à une position.\protect\hfil\protect\break}\label{lst:iv3:rd:lostmerge}
\end{question}

\begin{question}[$\star\star$ \iqroles{QD, Dev} \iqfirm{banque} \iqlevel{junior, confirmé}]\label{iq:iv3:reading-debugging-and-software-design:25}
\Cref{lst:iv3:rd:eventleak} est compilé en mode Release sur .NET 8. La \texttt{Strategy} que crée \texttt{Make} sort de portée, et rien dans \texttt{Main} ne la référence. Qu'affiche le programme après la collecte, et pourquoi~? Pourquoi de nombreuses stratégies de courte durée de vie poseraient-elles problème ici~? Donnez la correction habituelle.
\omcode{code/interviews-3/09-reading-debugging-and-software-design/csharp/EventLeak/Program.cs}{1}{29}{Un abonné et une collecte de ramasse-miettes.\protect\hfil\protect\break}\label{lst:iv3:rd:eventleak}
\end{question}

\begin{question}[$\star\star$ \iqroles{QR, QD, Dev, DS} \iqfirm{toutes} \iqlevel{junior, confirmé}]\label{iq:iv3:reading-debugging-and-software-design:26}
Les trois tests du \cref{lst:iv3:rd:weaktests} réussissent sous pytest. Combien d'entre eux pourraient jamais échouer~? Qu'a de faux \texttt{vwap}, et chacun des tests~? Réécrivez les tests de façon à détecter le bogue.
\omcode{code/interviews-3/09-reading-debugging-and-software-design/python/rd_vwap_weak.py}{2}{21}{Une fonction et ses tests.\protect\hfil\protect\break}\label{lst:iv3:rd:weaktests}
\end{question}

\begin{question}[$\star\star$ \iqroles{QD, Dev, LLE} \iqfirm{teneur de marché, hedge fund} \iqlevel{junior, confirmé}]\label{iq:iv3:reading-debugging-and-software-design:27}
\Cref{lst:iv3:rd:swapped} envoie 101 lots à un prix de 500, et le compilateur ne dit rien. Concevez les types de sorte que cette erreur ne se compile pas. Un collègue craint que les enveloppes coûtent du temps dans le chemin critique (hot path). Avec g++ 11.4, \texttt{-std=c++20 -O2}, sur x86-64 Linux, à quoi vous attendez-vous dans le code machine~?
\omcode{code/interviews-3/09-reading-debugging-and-software-design/cpp/swapped_args.cpp}{3}{10}{Un ordre avec deux longs.\protect\hfil\protect\break}\label{lst:iv3:rd:swapped}
\end{question}

\begin{question}[$\star\star$ \iqroles{QD, Dev, QR} \iqfirm{toutes} \iqlevel{junior, confirmé}]\label{iq:iv3:reading-debugging-and-software-design:28}
Un test a commencé à échouer quelque part dans les 1024 derniers commits. Chaque exécution, construction comprise, dure 6 minutes. Combien d'exécutions une recherche dichotomique demande-t-elle, et combien de temps dure-t-elle~? Maintenant, le test est instable~: sur un commit contenant le bogue, il échoue 1 fois sur 4, et sur un commit sans le bogue, il n'échoue jamais. Combien d'exécutions par commit faut-il pour que la recherche complète trouve le bon commit avec une probabilité d'au moins 99~\%~? Combien cela dure-t-il dans le pire des cas~?
\end{question}

\begin{question}[$\star\star$ \iqroles{QD, Dev} \iqfirm{banque} \iqlevel{junior, confirmé}]\label{iq:iv3:reading-debugging-and-software-design:29}
Vous relisez le \cref{lst:iv3:rd:readonly}. Un livre reconstruit à partir de l'historique ne doit pas recevoir de nouveaux ordres~; son auteur en a fait une sous-classe qui lève une exception. Le code s'exécute. Que dites-vous dans la relecture, et quelle forme demanderiez-vous à la place~? Existe-t-il un cas où vous l'accepteriez~?
\omcode{code/interviews-3/09-reading-debugging-and-software-design/python/rd_read_only.py}{1}{21}{Un livre en lecture seule.\protect\hfil\protect\break}\label{lst:iv3:rd:readonly}
\end{question}

\begin{question}[$\star\star\star$ \iqroles{QD, Dev, LLE} \iqfirm{toutes} \iqlevel{confirmé}]\label{iq:iv3:reading-debugging-and-software-design:30}
\Cref{lst:iv3:rd:uninit} affiche 200 dans chaque test unitaire et affiche parfois n'importe quoi en production. Compilé avec g++ 11.4, \texttt{-O2 -Wall -Wextra -Werror}, il ne produit aucun avertissement. Quel est le bogue, et pourquoi le test réussit-il~? Comment feriez-vous pour que la défaillance se produise à chaque fois, sur-le-champ, sans outil nouveau et sans recompilation~? Qu'affiche-t-elle alors~? Donnez la correction en un seul jeton.
\omcode{code/interviews-3/09-reading-debugging-and-software-design/cpp/uninit.cpp}{3}{21}{Une moyenne avec un bloc de statistiques.\protect\hfil\protect\break}\label{lst:iv3:rd:uninit}
\end{question}

\begin{question}[$\star\star\star$ \iqroles{QD, Dev, LLE} \iqfirm{teneur de marché, hedge fund} \iqlevel{confirmé}]\label{iq:iv3:reading-debugging-and-software-design:31}
La mémoire résidente d'un processus de trading augmente de quelques centaines de mégaoctets par jour. Vous l'exécutez sous AddressSanitizer avec la détection de fuites activée, et à la sortie il ne signale rien. Le \Cref{lst:iv3:rd:growth} montre la forme du code. Pourquoi l'outil reste-t-il silencieux~? Trois choses différentes peuvent faire croître la mémoire. Nommez-les, et décrivez comment vous identifiez laquelle vous avez, dans l'ordre.
\omcode{code/interviews-3/09-reading-debugging-and-software-design/cpp/growth.cpp}{4}{13}{Des ordres conservés pour l'audit.\protect\hfil\protect\break}\label{lst:iv3:rd:growth}
\end{question}

\begin{question}[$\star\star\star$ \iqroles{QD, Dev, LLE} \iqfirm{toutes} \iqlevel{confirmé}]\label{iq:iv3:reading-debugging-and-software-design:32}
\Cref{lst:iv3:rd:dcl} semble soigneux~: il vérifie, verrouille, puis vérifie à nouveau. Compilé avec g++ 11.4 pour x86-64, est-il correct~? Que peut-il exactement se passer de mal, et le modèle mémoire fort de x86-64 le sauve-t-il~? Donnez la version correcte la plus simple, puis une version qui conserve la double vérification. Comment ce schéma se corrige-t-il en Java~?
\omcode{code/interviews-3/09-reading-debugging-and-software-design/cpp/dcl.cpp}{6}{24}{Verrouillage à double vérification.\protect\hfil\protect\break}\label{lst:iv3:rd:dcl}
\end{question}

\begin{question}[$\star\star\star$ \iqroles{QD, Dev, LLE} \iqfirm{toutes} \iqlevel{confirmé}]\label{iq:iv3:reading-debugging-and-software-design:33}
Un programme tient en trois fichiers. Le \Cref{lst:iv3:rd:rates} et \cref{lst:iv3:rd:limits} définissent chacun une variable globale, et \cref{lst:iv3:rd:simain} affiche \texttt{limit}, alors que \texttt{BASE\_RATE} n'est pas défini. Un collègue modifie l'ordre des fichiers objets dans la ligne d'édition de liens, et la limite de risque devient nulle. Expliquez cela. Que s'affiche exactement dans chaque ordre avec g++ 11.4 et GNU ld~? Donnez une correction qui fonctionne dans n'importe quel ordre, et un mot-clé qui transforme l'erreur en erreur de compilation.
\omcode{code/interviews-3/09-reading-debugging-and-software-design/cpp/static_rates.cpp}{3}{8}{Le fichier des taux.\protect\hfil\protect\break}\label{lst:iv3:rd:rates}
\omcode{code/interviews-3/09-reading-debugging-and-software-design/cpp/static_limits.cpp}{1}{3}{Le fichier des limites.\protect\hfil\protect\break}\label{lst:iv3:rd:limits}
\omcode{code/interviews-3/09-reading-debugging-and-software-design/cpp/static_main.cpp}{3}{5}{Le fichier principal.\protect\hfil\protect\break}\label{lst:iv3:rd:simain}
\end{question}

\begin{question}[$\star\star\star$ \iqroles{LLE, QD, Dev} \iqfirm{teneur de marché, hedge fund} \iqlevel{confirmé}]\label{iq:iv3:reading-debugging-and-software-design:34}
\Cref{lst:iv3:rd:dispatch} présente cinq façons d'appeler \texttt{score}. Compilé avec g++ 11.4 et sa bibliothèque standard, \texttt{-std=c++20 -O2}, sur x86-64~: lesquelles des cinq fonctions contiennent un appel ou un saut indirect~? Laquelle est la surprise~? Au-delà de la performance, qu'est-ce qui tranche entre elles lorsque vous concevez un moteur de stratégies~?
\omcode{code/interviews-3/09-reading-debugging-and-software-design/cpp/dispatch.cpp}{3}{25}{Cinq formes de répartition.\protect\hfil\protect\break}\label{lst:iv3:rd:dispatch}
\end{question}

\begin{question}[$\star\star\star$ \iqroles{QD, Dev, QT} \iqfirm{teneur de marché, hedge fund, crypto} \iqlevel{confirmé}]\label{iq:iv3:reading-debugging-and-software-design:35}
\Cref{lst:iv3:rd:replay} modélise ce que fait un gestionnaire de positions lors d'une reconnexion, lorsque la passerelle renvoie la dernière exécution. Qu'affiche-t-il, et que devrait-il afficher~? Repensez \texttt{on\_fill} de sorte qu'un renvoi soit sans effet. Que doit faire votre conception lorsqu'une exécution manque, au lieu d'être répétée~? Comment la testez-vous~?
\omcode{code/interviews-3/09-reading-debugging-and-software-design/python/rd_replay.py}{1}{15}{Une position lors d'une reconnexion.\protect\hfil\protect\break}\label{lst:iv3:rd:replay}
\end{question}

\begin{question}[$\star\star\star$ \iqroles{QD, Dev, QR, QT} \iqfirm{hedge fund, teneur de marché} \iqlevel{confirmé}]\label{iq:iv3:reading-debugging-and-software-design:36}
Vous voulez que la stratégie du \cref{lst:iv3:rd:strategy} s'exécute sans modification en trading réel et en backtest. Vous rejouez cinq ticks, espacés de 2 secondes en temps de marché, aussi vite que la machine le permet, tous au-dessus du niveau. Combien d'ordres la stratégie envoie-t-elle~? Combien devrait-elle en envoyer~? Repensez-la. Qu'est-ce qui, au-delà de l'horloge, empêche un backtest de se reproduire exactement à l'identique~?
\omcode{code/interviews-3/09-reading-debugging-and-software-design/python/rd_strategy.py}{1}{15}{Une stratégie limitée en débit.\protect\hfil\protect\break}\label{lst:iv3:rd:strategy}
\end{question}

\begin{question}[$\star\star\star$ \iqroles{QD, Dev, QR} \iqfirm{teneur de marché, hedge fund} \iqlevel{confirmé}]\label{iq:iv3:reading-debugging-and-software-design:37}
Les deux tests du \cref{lst:iv3:rd:propbook} réussissent, et \texttt{Book} contient encore un bogue. Sans chercher ce bogue à la lecture du code, écrivez un test aléatoire qui le trouve~: que génère-t-il, que vérifie-t-il après chaque étape, et comment obtient-il un cas court qui échoue plutôt qu'une liste de trente opérations~? À quoi ressemble le cas court~? Quel est le bogue~?
\omcode{code/interviews-3/09-reading-debugging-and-software-design/python/rd_prop_book.py}{1}{30}{Un petit livre et ses tests unitaires.\protect\hfil\protect\break}\label{lst:iv3:rd:propbook}
\end{question}

\begin{question}[$\star\star\star$ \iqroles{Dev, QD, LLE} \iqfirm{teneur de marché, hedge fund, banque} \iqlevel{confirmé}]\label{iq:iv3:reading-debugging-and-software-design:38}
Un message binaire de votre bus interne est lu par six services et écrit par deux. Vous devez ajouter un champ et en retirer un autre, sans pouvoir tout arrêter d'un coup. Comment déployez-vous le changement, et dans quel ordre~? Quelles règles le format doit-il respecter, et quels tests ajoutez-vous~? Vous conservez aussi des années de messages enregistrés~: qu'est-ce que cela change~?
\end{question}

\begin{question}[$\star\star\star\star$ \iqroles{QD, Dev, LLE} \iqfirm{teneur de marché, hedge fund} \iqlevel{confirmé}]\label{iq:iv3:reading-debugging-and-software-design:39}
\Cref{lst:iv3:rd:bookbugs} est un côté d'un carnet d'ordres, écrit pour un test à faire chez soi. Il sert pour le côté acheteur et pour le côté vendeur. Vous disposez de dix minutes pour le relire. Il y a au moins trois bogues. Pour chacun, donnez la plus petite suite d'appels qui le montre, et dites ce que vous lanceriez pour le voir. Lequel d'entre eux AddressSanitizer ne détecte-t-il pas~? Dites ensuite ce que vous changeriez dans la conception.
\omcode{code/interviews-3/09-reading-debugging-and-software-design/cpp/book_bugs.hpp}{7}{46}{Un côté d'un livre.\protect\hfil\protect\break}\label{lst:iv3:rd:bookbugs}
\end{question}

\begin{question}[$\star\star\star\star$ \iqroles{QD, Dev, QR, DS} \iqfirm{toutes} \iqlevel{junior, confirmé}]\label{iq:iv3:reading-debugging-and-software-design:40}
La réponse d'un candidat au test à faire chez soi est le \cref{lst:iv3:rd:review}~: elle lit un CSV de transactions aux colonnes \texttt{symbol}, \texttt{price} et \texttt{qty}, et calcule le prix moyen pondéré par les volumes par symbole. Elle donne les bons chiffres sur le fichier d'exemple. Vous la relisez. Listez ce qui ne va pas, du plus grave au moins grave, et pour chaque point dites quelle entrée le révèle. Esquissez ensuite la version que vous accepteriez, et les tests qui vont avec.
\omcode{code/interviews-3/09-reading-debugging-and-software-design/python/rd_review.py}{2}{26}{La réponse à un test à faire chez soi.\protect\hfil\protect\break}\label{lst:iv3:rd:review}
\end{question}
